\documentclass[a4paper]{article}
% Español (México) con XeLaTeX: fontspec para fuentes Unicode, polyglossia
% para la división silábica y los nombres del documento en la variante mexicana.
\usepackage{fontspec}
\usepackage{polyglossia}
\setdefaultlanguage[variant=mexican]{spanish}
% Los nombres chinos van como «pinyin (original)»: xeCJK da a los caracteres
% Han su propia fuente; sin ella XeLaTeX los omite con un aviso.
% FandolSong no cubre algunos caracteres de nombres (U+73FA); WenQuanYi Zen Hei sí.
\usepackage[AutoFallBack=true]{xeCJK}
\setCJKmainfont{FandolSong-Regular.otf}
\setCJKfallbackfamilyfont{\CJKrmdefault}{WenQuanYi Zen Hei}
% polyglossia no traduce los nombres de listings.
\AtBeginDocument{\providecommand{\lstlistingname}{}\renewcommand{\lstlistingname}{Listado}%
  \providecommand{\lstlistlistingname}{}\renewcommand{\lstlistlistingname}{Índice de listados}}
\usepackage{amsmath, amssymb}
\usepackage{graphicx}
\usepackage[margin=2.5cm]{geometry}
\usepackage[most]{tcolorbox}
\usepackage{etoolbox}
\usepackage{listings}
\usepackage{hyperref}
\usepackage{booktabs}
\usepackage{subcaption}
\usepackage{float}

\newtcolorbox{knowledgebox}[1]{enhanced, colback=blue!5!white, colframe=blue!75!black, colbacktitle=blue!75!black, coltitle=white, fonttitle=\bfseries, title={#1}, attach boxed title to top left={yshift=-2mm, xshift=2mm}, boxrule=1pt, sharp corners}
\newtcolorbox{importantbox}[1]{enhanced, colback=yellow!10!white, colframe=yellow!80!black, colbacktitle=yellow!80!black, coltitle=black, fonttitle=\bfseries, title={#1}, sharp corners}
\newtcolorbox{warningbox}[1]{enhanced, colback=red!5!white, colframe=red!75!black, colbacktitle=red!75!black, coltitle=white, fonttitle=\bfseries, title={#1}, sharp corners}

\lstset{language=Python, basicstyle=\ttfamily\small, keywordstyle=\color{blue}, stringstyle=\color{red!60!black}, commentstyle=\color{green!60!black}, breaklines=true, frame=single, numbers=left, numberstyle=\tiny\color{gray}, captionpos=b, extendedchars=true}

\newcommand{\notetitle}{[LLM Agents SP25] AlphaProof: RL Meets Formal Mathematics — Thomas Hubert}
\newcommand{\noteauthors}{Elaboradas a partir de la clase de Thomas Hubert}
\newcommand{\notedate}{10 de marzo de 2025}
\newcommand{\videochannel}{Berkeley RDI}
\newcommand{\videopublishdate}{7 de marzo de 2025}
\newcommand{\videoduration}{aproximadamente 60 minutos}
\newcommand{\videourl}{}
\newcommand{\videocoverpath}{cover.jpg}
\newcommand{\repourl}{https://github.com/hqhq1025/ai-course-notes}


% Footer with repo info
\usepackage{fancyhdr}
\pagestyle{fancy}
\fancyhf{}
\fancyfoot[C]{\thepage}
\fancyfoot[R]{\footnotesize\color{gray}\href{\repourl}{AI Course Notes}}
\renewcommand{\headrulewidth}{0pt}
\renewcommand{\footrulewidth}{0pt}

\begin{document}
\begin{titlepage}
\centering
{\Large Notas del curso\par}\vspace{1.2cm}
{\huge\bfseries \notetitle\par}\vspace{0.8cm}
{\large \noteauthors\par}\vspace{0.3cm}
{\large \notedate\par}\vspace{1.2cm}
\ifdefempty{\videocoverpath}{}{\includegraphics[width=0.82\textwidth,height=0.45\textheight,keepaspectratio]{\videocoverpath}\par}
\vfill
\begin{tcolorbox}[width=0.9\textwidth, colback=black!2!white, colframe=black!60, sharp corners]
\textbf{Autor o canal del video}: \videochannel\par
\textbf{Fecha de publicación}: \videopublishdate\par
\textbf{Duración del video}: \videoduration\par
\ifdefempty{\videourl}{}{\textbf{Enlace del video}: \href{\videourl}{\nolinkurl{\videourl}}\par}\par
\textbf{Página del proyecto}: \href{\repourl}{\nolinkurl{\repourl}}\par
\end{tcolorbox}
\end{titlepage}
\tableofcontents
\newpage

\section{Introducción: de AlphaGo a AlphaProof}
Thomas Hubert ubica esta clase en la unidad de «razonamiento-ejecución» de Berkeley LLM Agents SP25, y usa AlphaZero/AlphaTensor para enlazar la época dorada del aprendizaje por refuerzo con AlphaProof. Hubert señala explícitamente que la \textbf{rigurosidad}, la \textbf{feedback demostrable} y la \textbf{creatividad ilimitada} que exige el razonamiento matemático son altamente isomorfas a la \textbf{búsqueda}, la \textbf{generalización} y el \textbf{retorno confiable} que exige la exploración de fronteras en RL. Esto explica por qué DeepMind sitúa a las matemáticas como «el siguiente campo de batalla de RL».

\begin{knowledgebox}{La resonancia entre las matemáticas y RL}
Las matemáticas son una red de alta dimensión donde los nodos son teoremas y las aristas son cadenas lógicas; RL, por su parte, explora en el espacio de acciones. Hubert enfatiza: «en cuanto puedas abstraer un teorema como una política de acciones, el optimizador puede encargarse de aprender la teoría.»
\end{knowledgebox}

La clase se divide en tres partes: la primera repasa la evolución de la formalización matemática y organiza Lean/Mathlib; la segunda resume el ciclo virtuoso de AlphaZero y cómo AlphaProof se conecta con los sistemas de demostración; la tercera se centra en la práctica de ingeniería, la reproducción de la IMO y el ciclo de gobernanza. Toda la clase sigue la lógica pedagógica de «teoría → plataforma → aplicación», y en cada parte se reserva tiempo para mostrar lecciones aprendidas y preguntas y respuestas.

\begin{figure}[H]
\centering
\includegraphics[width=0.8\textwidth,height=0.35\textheight,keepaspectratio]{\videocoverpath}
\caption{Portada de la Lecture 05 y esquema del flujo de AlphaProof (la portada aporta la imagen principal y una imagen arquitectónica, y resalta la fusión de RL con la formalización).}
\end{figure}

Hubert menciona en particular que esta clase es la quinta de SP25, e inmediatamente después de «agentes generativos y razonamiento». Considera a AlphaProof como «un agente en la versión matemática», que necesita combinar el retorno determinista con la gobernanza interactiva para que la búsqueda de demostraciones de alta confiabilidad sea posible.

\subsection{Resumen de la sección}
La parte introductoria se encarga de construir el contexto: la similitud estructural entre las matemáticas y RL, la misión de AlphaProof, y la ubicación en el ritmo de toda la serie de clases de SP25 entre la «cadena de razonamiento y la práctica de ingeniería», sentando así la ruta causal para el contenido posterior.
\newpage

\section{Historia y ecosistema de la formalización matemática}

\subsection{De Grecia a la computadora: generaciones de la formalización}
La demostración matemática comenzó en Grecia; Euclides inauguró la axiomatización e Hilbert llevó la formalización al lenguaje. De la simbolización a la era del cómputo, las matemáticas se apoyaron cada vez más en la verificación automatizada: la demostración por inducción de Gentzen, la correspondencia de Curry-Howard, Coq e Isabelle se acercan a la intersección entre «legible por humanos» y «verificable por máquina».
\begin{itemize}
    \item 500 a.C.: la geometría griega traza la red axiom-theorem;
    \item 1900s: el programa de Hilbert y el sistema axiomático formal;
    \item 1970s: Curry-Howard hace que el tipo sea la demostración;
    \item 2010s: Lean y Mathlib adoptan la colaboración abierta, con la meta de que «cualquiera pueda formalize»;
    \item 2020s: AlphaProof conecta un motor de búsqueda de RL a Lean, para abordar proposiciones de nivel IMO.
\end{itemize}

Hubert también subraya que la formalización matemática no es una historia unidireccional, sino una «actualización parallel de múltiples generaciones»: cada PR de Mathlib se replay en el entorno de RL, y el RL agent a su vez retroalimenta a Mathlib con el trace de sus fallas, mientras que los governers forman discusiones interpretables en GitHub/CI; esta retroalimentación cruzada impulsa la adopción de nuevos axiomas.

\subsection{Lean, Mathlib y la estandarización}
\begin{importantbox}{Los tres principios de Lean + Mathlib}
1) \textbf{Reproducibilidad}: cada teorema cuenta con un script compilable; 2) \textbf{Unificación}: cada concepto entra en una sola biblioteca, para evitar la fragmentación de definiciones; 3) \textbf{Apertura}: Mathlib convierte a los colaboradores en nodos de gobernanza mediante PR/CI. Hubert subraya: el RL agent debe entender Mathlib y también ser sensible a los fallos de CI.
\end{importantbox}

Estos principios se reflejan directamente en el despliegue de AlphaProof: el formalizer se alinea con el pipeline de PR/CI de Mathlib, y cada vez que el agent genera un nuevo lemma debe pasar primero por el nightly runner; cuando el CI failed, el agent degrada automáticamente su search policy y, en su lugar, comparte el trace con el grupo human-in-the-loop.

Elementos semánticos que Lean ofrece a AlphaProof:
\begin{enumerate}
    \item \textbf{state del entorno}: el teorema actual, las premisas y la pila de tactic;
    \item \textbf{catalog de acciones}: los distintos tactic conforman el action space, que puede personalizarse con BFS, MCTS o Heuristic;
    \item \textbf{Retroalimentación}: la demostración exitosa equivale a +1, y una failure trigger un trace/log;
    \item \textbf{Observation}: Lean regresa el proof term, el árbol de goal y las restricciones de type.
\end{enumerate}

\subsection{Mecanismos de colaboración en el ecosistema de formalización}
Hubert describe el ciclo cerrado de colaboración de la comunidad de Lean: los mantenedores de Mathlib controlan la calidad mediante el pipeline de PR/build, y el equipo de AlphaProof sube «proof search logs + failed tactics» al governance board, lo que forma la entrada del trace de RL. Varios Lean expert abren el replay log para corregir hallucination de manera human-in-the-loop.

\begin{knowledgebox}{Puntos clave de la formalización colaborativa}
1) Crowdsourced replay log: registra la sequence de fallas del agent; 2) Fairness dashboard: verifica si la search está sesgada hacia algún tipo de política; 3) Human review loop: cuando el RL activa el drift gate, el expert puede rollback la knowledge base.
\end{knowledgebox}

\subsection{Resumen de la sección}
El desarrollo de la formalización matemática, los principios de estandarización de Lean/Mathlib y los mecanismos de colaboración community-driven sostienen en conjunto el contexto de AlphaProof; esta sección explica por qué conectar RL con Lean produce «creatividad verificable».
\newpage

\section{Logros y umbrales del aprendizaje por refuerzo}

\subsection{Filosofía y método de la serie Zero}
Hubert atribuye el éxito de AlphaZero a dos ejes centrales: \textbf{search-integrated learning} y \textbf{reliable reward from simulation}. En los juegos, el entorno simulado ofrece perfect feedback, de modo que el RL ya no depende de etiquetas humanas, sino que usa las victorias y derrotas de la búsqueda como verdad de referencia.
\begin{knowledgebox}{Los tres pilares de la filosofía Zero}
1) Self-play generating curriculum; 2) Search results distilled into policy/value; 3) Preference for small, reusable modules rather than monolithic black-box policies.
\end{knowledgebox}

Él dijo: «El esplendor de la serie Zero no es el éxito de un solo modelo, sino de un proceso de ingeniería depurado con precisión: cada self-play produce un curriculum, y la policy/value se reconstruye como un heuristic module listo para la búsqueda». Esta filosofía influye directamente en las expectativas de policy distillation de AlphaProof.

\subsection{El volante de búsqueda-aprendizaje}
AlphaZero/AlphaTensor construyeron un volante de «exploración → aprendizaje → optimización». Hubert señala en particular: sin strong search no se obtienen high-quality trajectories; sin fast learning no es posible interiorizar el conocimiento nuevo; y sin una estrecha combinación entre policy y value, la search cae en un thousand-node dead-end.
\begin{warningbox}{Vacíos del volante y fallas potenciales}
1) Search de baja calidad: el agent no logra encontrar soluciones sólidas y la actualización de policy es débil; 2) reward con ruido: la etapa de aprendizaje incorpora hallucinated trajectories; 3) Compute limitado: una search depth insuficiente impide la convergence del learning.
\end{warningbox}

Hubert compara este volante con «excavar junto a un agujero negro»: la search no puede quedarse en la superficie, o el learning nunca tendrá con qué avanzar; el learning también debe absorb a tiempo los nuevos lemma que descubre la search, o la search seguirá repitiendo el mismo camino. El equipo de AlphaProof usa un two-tier policy buffer para que search-area y learning-rate se sigan mutuamente y evitar así el drift.

\subsection{Exploración, retroalimentación y generalización}
La retroalimentación de AlphaZero es un win-rate determinístico; AlphaProof, en cambio, debe enfrentar la «corrección formal de la demostración»: si la prueba en Lean no pasa, la retroalimentación es simplemente un fallo. Hubert explica: «En matemáticas, la reward es +1 o es 0, no hay zona gris, y eso obliga al RL a desarrollar trace-aware regret minimization». Comparación entre búsqueda y heurísticas:
\begin{table}[H]
\centering
\begin{tabular}{p{4cm}p{4cm}p{4cm}}
\toprule
Dimensión & Search-driven RL & Heuristic/Scripted \\
\midrule
Calidad de la retroalimentación & Verificación perfecta, replay buffer de alta precisión & Señal con mucho ruido, propensa al drift \\
Capacidad de generalización & Logra la transferencia entre teoremas mediante policy distillation & Requiere ingeniería manual \\
Capacidad de auditoría & Se puede reproducir la trajectory y ofrecer el proof term & Difícil de reducir a pasos verificables \\
\bottomrule
\end{tabular}
\caption{Comparación entre la retroalimentación de RL y los métodos heurísticos en AlphaProof}
\end{table}

\subsection{Resumen de la sección}
Los tres planos del éxito de AlphaZero (search, feedback, learning) son el punto de partida de AlphaProof; solo prestando atención a cualquier ruptura del volante, al colapso de la señal de retroalimentación y a los límites de la generalización se puede asegurar que el RL agent avance de forma estable en el mundo matemático.
\newpage

\section{AlphaProof: RL se encuentra con las matemáticas formales}

\subsection{Lean como el entorno de RL ideal}
Lean ofrece una tupla completa de state/action/reward/observation. Hubert mostró un episode típico: a partir del `goal` se abstrae un lemma graph → el agent elige `apply`/`intro`/`rw` → Lean intenta simplificar el goal → éxito o fail. «Al final tratamos la salida de Lean como el simulator del entorno de RL», dijo.
\begin{importantbox}{Anatomía del entorno}
1) Action space: tactics en capas (primitive tactic, combination tactic, search macro); 2) Observation: árbol de goal + local context + type info; 3) Reward: binario; 4) Reset: en cada new theorem se reinicializa el árbol de MCTS.
\end{importantbox}

\subsection{Arquitectura del sistema y pipeline de entrenamiento}
AlphaProof está compuesto por varias redes: el Formalizer traduce proposiciones en lenguaje natural/LaTeX en aserciones de Lean; el Proof Generator/Actor-critic se encarga de la search; el Verifier realiza una segunda verificación; el Runner registra el trace.
\begin{table}[H]
\centering
\begin{tabular}{p{4.5cm}p{4.5cm}p{4cm}}
\toprule
Componente & Responsabilidad & Notas \\
\midrule
Formalizer & Estructura las proposiciones, genera declarations de Lean & Basado en seq-to-seq + retrieval \\
Proof Generator & RL + search genera tactic sequence & Combinación de Monte Carlo y beam \\
Verifier & Lean verifica el proof term & Si falla, dispara un rollback \\
Runner & Registra proof trace + metrics & Se usa para el governance gate \\
\bottomrule
\end{tabular}
\caption{Componentes de la arquitectura de AlphaProof}
\end{table}
\begin{knowledgebox}{Puntos de gobernanza del pipeline}
Cada component escribe en el governance log: la comparación de semantic gap del Formalizer, la entropy/confidence del Proof Generator, la failure rate del Verifier, el chunk gating hit del Runner. Hubert dijo: «Solo si convertimos el trace en un artifact estructurado, el RL puede dejar de ser una caja negra».
\end{knowledgebox}

Durante el entrenamiento, AlphaProof alterna el entrenamiento entre varios dataset: Mathlib core definitions, IMO contest statements, AlphaGeometry proofs. Hubert enfatizó que cada dataset corresponde a un replay buffer independiente (el «Lemma buffer»), y que en el training schedule se establecen 3 etapas, warm start → search-intensive → human eval, lo que hace que la policy sea cada vez más robusta.

\subsection{Experimentos, IMO 2024 e iteración}
AlphaProof, junto con AlphaGeometry, obtuvo la medalla de plata en el IMO 2024; los principales avances fueron: 1) construir un large-scale formalization dataset; 2) usar el replay buffer para registrar el failure trace durante la proof search; 3) usar el resultado del human review del governance board como reward shaping. Hubert mostró en particular un ejemplo: cuando el agent tiene un fail en el teorema `prime`, el trace se sube a `AlphaProof Repos`, un math expert genera un corrected lemma, y este se feed de nuevo a la policy.
\begin{warningbox}{Advertencias para el despliegue a nivel IMO}
1) hay que probar de forma segmentada parámetros como search depth, tactic mix, chunk gating; 2) cuando el tiempo de permanencia del proof supera los 500 ms hay que forzar el log; 3) no basta con mirar el final success rate, también hay que hacer tally de los drift gate triggers.
\end{warningbox}

\subsection{Resumen de la sección}
AlphaProof aprovecha el deterministic reward, structured observation y rich action space que ofrece Lean para convertir la math proof search en un RL pipeline controlable; el log de gobernanza y el human review hacen que incluso el fracaso sea útil.
\newpage

\section{Evaluación y modos de falla}

\subsection{Jerarquía de métricas y gobernanza}
\begin{table}[H]
\centering
\begin{tabular}{p{3.5cm}p{4cm}p{3.5cm}p{3cm}}
\toprule
Métrica & Descripción & Frecuencia de recolección & Condición de activación \\
\midrule
Proof success rate & Si Lean demuestra con éxito la proposición & real-time per episode & success < 2\% → ampliar search \\
Entropy drift & Policy entropy drop & epoch & drop > 30\% → drift gate log \\
Chunk hit rate & Tasa de aciertos de la cache del search track & per batch & < 80\% → adjust overlap \\
Human edit count & Número de traces que requieren intervención humana & weekly & > 5/1000 → governance review \\
\bottomrule
\end{tabular}
\caption{Métricas de evaluación multinivel y estrategias de respuesta de AlphaProof}
\end{table}

Hubert mencionó que las diapositivas 17-19 muestran el dashboard completo de métricas: la parte superior corresponde a proof success/distribution, el lado izquierdo al entropy drift heatmap y el lado derecho al user feedback count. El dashboard está conectado directamente con el alert pipeline del governance board, y en cuanto una track cruza el threshold, automáticamente envía el replay log al reviewer.

\begin{importantbox}{El ciclo cerrado de evaluación y gobernanza}
1) Recolección paralela de múltiples métricas; 2) cuando el proof success rate cae por debajo del baseline, se aumenta el search depth y se registra el replay; 3) el entropy drift y los manual edits determinan en conjunto si se necesita una governance review.
\end{importantbox}

\subsection{Análisis de los modos de falla}
Las fallas registradas por AlphaProof se concentran sobre todo en los siguientes patterns: A) semantic gap: error de formalizer translation; B) search drift: la policy entra en una low-entropy zone; C) Lean constraints: un type mismatch provoca el colapso del goal tree; D) compute budget exceed. Hubert lo representa de manera gráfica con el «failure funnel»: la mayoría de las failure se filtran en la etapa de initial seeding, y solo una minoría hace drift via chunk gating hasta llegar al governance board.

\begin{warningbox}{Advertencia sobre fallas comunes}
El semantic gap requiere volver al formalizer; una vez que se activa el drift gate, primero se analiza el entropy heatmap; los errores de tipo requieren revisar el Lean goal stack; cuando ocurre un budget exceed, se baja directamente al low-latency track.
\end{warningbox}

\subsection{Reproducción y robustez}
Cada failure trace se escribe en `failure-repro.json`, e incluye time stamp, search track, entropy y human comment. El equipo sube estos trace al experiment repo para generar el `failure curriculum`, que permite que el RL agent haga un nuevo warm start. Este proceso está supervisado por el governance board, que garantiza que cada failure log pueda convertirse en el next training seed.

\begin{knowledgebox}{Elementos del failure curriculum}
1) Trace + entropy + comment conforman una entry; 2) la entry entra en el Lemma buffer; 3) el training schedule la usa como positive/negative sample; 4) Result logged back to dashboard.
\end{knowledgebox}

\subsection{Resumen de la sección}
La capa de evaluación y modos de falla conecta la percepción (métricas), el diagnóstico (failure taxonomy) y la reproducción (failure curriculum), lo que permite que AlphaProof, a lo largo de sus iteraciones, convierta las fallas en datos de entrenamiento controlables.
\newpage

\section{Caso práctico: búsqueda multiestrategia del teorema de números primos}

\subsection{Descomposición del problema y semantización}
Hubert descompuso un teorema de números primos en tres subgoals: 1) construir candidatos numéricos a partir de la definición; 2) usar un lemma para vincular el número primo con el módulo; 3) usar una tactic para enumerar los steps de razonamiento. Cada subgoal forma un goal tree independiente en Lean; AlphaProof primero busca lemmas similares en Mathlib y luego invoca a la policy para generar la tactic.

\begin{knowledgebox}{Textura semántica del teorema de números primos}
1) capa de definición: el concepto de prime number necesita transitar de la aritmética básica a la teoría de números avanzada; 2) capa de estrategia: distintos lemmas como \texttt{mod\_pow} y \texttt{order\_dvd} usan tactics diferentes; 3) capa de gobernanza: cada tactic fallida se escribe en el trace, para que un human reviewer la revise después.
\end{knowledgebox}

\subsection{Combinación de estrategias y chunk gating}
En la ejecución real, AlphaProof mantiene a la vez tres search tracks: heuristic-run (bajo costo), beam-run (alta fidelidad) y MCTS-run (pulido profundo). La tabla de chunk gating registra, para cada track, el token budget, el overlap y el entropy drift; cuando heuristic-run esboza una posible solución, beam-run toma el relevo; si el drift supera el umbral, se degrada a MCTS-run, para garantizar la corrección de la eventual proof.

\begin{table}[H]
\centering
\begin{tabular}{p{3.5cm}p{3.5cm}p{5cm}}
\toprule
Track & Budget & Descripción \\
\midrule
Heuristic-run & 128 tokens & Exploración rápida que genera un proof sketch; cuando el heatmap baja de 0.4, deja entrar a beam \\
Beam-run & 256 tokens & Expansión conservadora de los top-k candidate, incorporando lemmas \\
MCTS-run & 512 tokens & Pulido profundo, activado por el drift gate \\
\bottomrule
\end{tabular}
\caption{Search tracks y gating del teorema de números primos}
\end{table}

\begin{warningbox}{Consideraciones sobre la planificación en tiempo real}
El chunk gating no puede ser solo una threshold table: también debe ajustar el budget según la longitud del proof term; si el drift gate se activa con frecuencia, indica que el parameter del Beam track es demasiado aggressive.
\end{warningbox}

\subsection{Resumen de la sección}
El caso del teorema de números primos muestra cómo la multi-track search se combina con el chunk gating para que la RL policy explore con rapidez sobre teoremas matemáticos sin perder precisión; el trace log convierte los fallos que encuentra el drift gate en insumo para la gobernanza.
\newpage
\section{Práctica de ingeniería y control}

\subsection{Planificación de Proof Search}
Hubert presentó un planificador: los Edges tienen distintos cost (Heuristic, Beam, MCTS), y en cada paso se elige la search policy en la tabla de chunk gating; chunk gating registra el token count, el Overlap y el proof term size. El planificador elige, según el latency budget, entre «search superficial + high replay» o «MCTS profundo + low replay».
\begin{table}[H]
\centering
\begin{tabular}{p{4cm}p{4cm}p{4cm}}
\toprule
Tipo de search & Características & Escenario de aplicación \\
\midrule
Heuristic tactics & Baja latency & exploration inicial \\
Beam search & selección conservadora de top-k & lemma discovery \\
MCTS + RL policy & value de estado + policy & polish final, adecuado para nivel de teorema \\
\bottomrule
\end{tabular}
\caption{Estrategia de planificación de Proof search}
\end{table}
\begin{importantbox}{KPI del planificador}
El latency budget, el proof-term length y el entropy drift determinan en conjunto el chunk gating gate; al superar el umbral, se degrada a conservative search.
\end{importantbox}

\subsection{Registro Trace y verificación}
Cada failed proof se registra en `trace.json`: incluye la tactic sequence, las indicaciones de Lean y el search depth. Hubert enfatizó que el trace debe tener un campo `human comment`, para que el governance team pueda hacer triage rápidamente. También señaló que, tras un fallo de verificación de Lean, la RL necesita replay de un alternate candidate para evitar quedar stuck.
\begin{warningbox}{Advertencia sobre el registro Trace}
Sin replay (trail) es imposible medir el strategy drift; un replay buffer sin trace hace que la policy entrene a ciegas con wrong proofs.
\end{warningbox}

\subsection{Gobernanza humano-máquina y reproducción}
El equipo de gobernanza envía periódicamente los high-drift chunk al human-in-the-loop board: el reviewer revisa el trace/log, evalúa la fairness y actualiza el knowledge summary en el newsletter. Esta práctica hace que AlphaProof trate cada fallo como seed de un experimento de reproducción.
\begin{knowledgebox}{Puntos clave de gobernanza}
1) High drift chunk activa un alert; 2) el reviewer genera un corrected tactic; 3) policy replay + reward correction; 4) el newsletter registra los puntos destacados de QA.
\end{knowledgebox}

\subsection{Resumen de la sección}
La parte de ingeniería destaca la planificación de Proof search, el registro de trace y el ciclo cerrado de gobernanza: mediante el triángulo de chunk gating, trace + review y newsletter, convierte al RL agent de una black box inexplicada en una línea vital de investigación y desarrollo reproducible.
\newpage

\section{Compartición de conocimiento e iteraciones futuras}

\subsection{Flujo de conocimiento y reproducción}
El equipo de AlphaProof registra el experiment log, el failed proof trace, los human comment y los experimentos del newsletter en un solo repo; cada release viene acompañado de un `diff` que se usa para reentrenar la policy. Hubert señala: «Necesitamos compartir con los nuevos integrantes los aspectos destacados de la QA, las governance decision y los proof de ejemplos positivos y negativos, para que el training loop pueda iterar más rápido».
\begin{importantbox}{Mecanismo de compartición de conocimiento}
1) Experiment repo, que incluye el chunk gating journal; 2) Newsletter, que registra cada semana los aspectos destacados de QA + governance; 3) Shared slides + transcript, que también son el material para el ramp-up de los nuevos integrantes.
\end{importantbox}

\subsection{Aspectos destacados de la QA}
La sesión de QA se usó para explicar si AlphaProof puede extenderse a la non-Euclidean geometry y al ~hypergraph theorem. Al responder, Hubert subrayó que la RL policy debe diseñarse de forma modular, ajustando cada vez solo el search depth o el tactic mix, para poder extenderse gradualmente a demostraciones más abstractas.
\begin{warningbox}{Puntos de precaución en la QA}
No se debe usar AlphaProof de manera indiscriminada para resolver todos los problemas; en los idiomas low-resource o en las novel math branch, la policy necesita un nuevo warm start, junto con el chunk gating y el fairness board.
\end{warningbox}

\subsection{Resumen de la sección}
El apartado sobre compartición de conocimiento mostró cómo el log, el newsletter y el transcript forman una red de reproducción, y cómo los retos de extensión expuestos mediante la sesión de QA explican mejor los límites de la iteration y las necesidades de gobernanza.
\newpage

\section{Implicaciones del sistema y gobernanza futura}

\subsection{Panel de implicaciones del sistema}
Hubert llamó al diseño de AlphaProof «arquitectura interpretable de tres capas»: 1) Formalizer/Proof Generator/Verifier conforman el deterministic pipeline; 2) Trace/log + dashboard constituyen la observability layer; 3) Governance review + newsletter son la human-in-the-loop layer. La experiencia de AlphaProof nos dice que un agent complejo debe separar con claridad en capas el pipeline, la metric y la gobernanza para poder desplegarse de forma segura en un ámbito sin matices de gris como math.

\begin{table}[H]
\centering
\begin{tabular}{p{3cm}p{4cm}p{5cm}}
\toprule
Nivel & Punto de atención & Método \\
\midrule
Pipeline & correctness + modularity & Formalizer/Proof Generator/Verifier + Runner conforman steeped pipeline \\
Observabilidad & metrics & Dashboard + chunk gating + trace log \\
Governance & trust & human review + newsletter + failure curriculum \\
\bottomrule
\end{tabular}
\caption{Las tres capas de implicaciones del sistema de AlphaProof}
\end{table}

\begin{knowledgebox}{Esencia de las implicaciones del sistema}
Concebir el system design como una arquitectura de tres capas — pipeline + observabilidad + governance — permite garantizar que la tarea de razonamiento matemático más sensible mantenga su interpretabilidad y su control incluso en escenarios desconocidos.
\end{knowledgebox}

\subsection{Dirección de la gobernanza futura}
Hubert propuso la idea de «differentiated governance»: fijar un drift threshold distinto para cada search track y registrar las governance decision en el newsletter/backlog. Por ejemplo, ante una vulnerabilidad de beam search, el governance reviewer aplica de manera directa un patch a la policy correspondiente, mientras que en el heuristic track es el RL agent quien reintenta de forma autónoma. Él enfatiza que la gobernanza no es un judge, sino una feedback policy.

\begin{warningbox}{La gobernanza no debe reducirse a prohibir}
Cuando se activa una governance alert, no se trata de desactivar directamente al agent, sino de escribir de vuelta el human comment en el failure curriculum, aumentar el reward shaping, y después reanudar la search; de lo contrario, el system terminará convertido en un fragile checkbox.
\end{warningbox}

\subsection{Resumen de la sección}
Las implicaciones del sistema se resumen en una arquitectura de tres capas, mientras que la gobernanza futura, mediante differential threshold + feedback loop, permite que el agent matemático siga expandiendo sus límites; estas reflexiones también pueden retroalimentar otros sistemas RL que requieren un alto grado de interpretabilidad.
\newpage

\section{Síntesis y ampliación}
\begin{table}[H]
\centering
\begin{tabular}{p{4cm}p{6cm}p{5cm}}
\toprule
Módulo & Aprendizaje clave & Siguiente paso \\
\midrule
Formalización matemática & Lean/Mathlib ofrece un deterministic environment y una gobernanza colaborativa & Incorporar el PR/gen replay log de Mathlib al experiment pipeline \\
Aprendizaje por refuerzo & El volante search-feedback-learning de AlphaZero es la base de diseño de AlphaProof & Continuar reforzando la search-policy distillation y el entropy gate \\
Arquitectura de AlphaProof & Formalizer + search + verifier + runner forman un pipeline controlable & Reforzar el mecanismo de activación de chunk gating y enriquecer el trace artifact \\
Evaluación y fallos & Las métricas multinivel y el failure curriculum convierten el failure en datos de entrenamiento & Ampliar el conjunto de métricas del dashboard y activar automáticamente el failure curriculum \\
Práctica de ingeniería & Proof search scheduling + trace logging + governance closing loop & Establecer el drift threshold y tratar cada alert como un retraining seed \\
Compartición de conocimiento & El log, el newsletter y el transcript forman un network de reproducibilidad & Hacer que el newsletter sea el archival record de los highlights de QA y de governance \\
Enseñanza del sistema & La arquitectura de tres capas de Pipeline/observability/governance constituye la base de la confianza & Incorporar la differential governance policy al newsletter \& la automation \\
\bottomrule
\end{tabular}
\caption{Módulos técnicos y plan de acción de la Lecture 05}
\end{table}

\subsection{Lecturas adicionales}
\begin{itemize}
    \item Hubert et al., ``AlphaProof: RL Meets Formal Mathematics,'' 2024.
    \item Documentación oficial de Lean y el README de Mathlib (GitHub).
    \item Artículos relacionados con AlphaZero / AlphaTensor.
    \item Archive público de silver proofs del IMO 2024.
    \item Blog de Terence Tao sobre la colaboración en formalización.
    \item Workshop notes de «Governance for autonomous agents».
\end{itemize}

\subsection{Resumen rápido de términos clave}
\begin{itemize}
    \item Chunk gating: selector de estrategia de token/overlap/latency dentro del proof search.
    \item Drift gate: la caída de entropy activa el fallback + review.
    \item Trace log: la tactic sequence + human comment de cada failed proof.
    \item Knowledge pipeline: el triángulo de log + newsletter + transcript.
\end{itemize}
\begin{knowledgebox}{El papel del resumen rápido de términos clave}
Usa una terminología unificada para ayudar a los nuevos integrantes a familiarizarse rápidamente con el operational playbook de AlphaProof, y se combina con el governance board para formar un ciclo cerrado de retroalimentación.
\end{knowledgebox}

\subsection{Resumen de la sección}
AlphaProof es un sistema que entrelaza la formalización matemática, el aprendizaje por refuerzo y una práctica de ingeniería gobernable: Lean ofrece el sandbox, RL ofrece el volante de search-policy, la gobernanza ofrece trace/alert/knowledge, y el newsletter mantiene sincronizado a todo el equipo.

\end{document}
